% ------------------------------------------------------------------
% Preámbulo común a la solución del taller y a la guía de estudio
% ------------------------------------------------------------------
\usepackage[T1]{fontenc}
\usepackage[utf8]{inputenc}
\usepackage[spanish,es-nodecimaldot,es-noshorthands,es-tabla,es-lcroman]{babel}
\usepackage{mathpazo}
\usepackage{tgpagella}
\usepackage[scaled=.92]{helvet}
\renewcommand{\ttdefault}{lmtt}
\linespread{1.06}
\usepackage{amsmath,amssymb,amsthm,mathtools}
\usepackage[letterpaper,margin=2.25cm,headheight=14pt]{geometry}
\usepackage[table,dvipsnames]{xcolor}
\usepackage{array,booktabs,longtable,multicol,graphicx}
\usepackage{enumitem}
\usepackage{microtype}
\usepackage{fancyhdr}
\usepackage{titlesec}
\usepackage{tikz}
\usetikzlibrary{arrows.meta,positioning,calc,shapes.geometric,fit,backgrounds}
\usepackage[most]{tcolorbox}
\usepackage[hidelinks,unicode]{hyperref}
\usepackage{bookmark}

% ----------------------------------------------------------- colores
\definecolor{tinta}{RGB}{24,52,94}
\definecolor{tinta2}{RGB}{46,96,150}
\definecolor{verde}{RGB}{22,102,82}
\definecolor{ambar}{RGB}{160,96,8}
\definecolor{rojo}{RGB}{160,36,36}
\definecolor{malo}{RGB}{250,214,214}
\definecolor{bueno}{RGB}{214,236,224}
\definecolor{bueno2}{RGB}{218,228,244}
\definecolor{papel}{RGB}{246,248,251}
\hypersetup{colorlinks=true,linkcolor=tinta2,urlcolor=tinta2,citecolor=tinta2}

% ---------------------------------------------------------- formato
\setlength{\parindent}{0pt}
\setlength{\parskip}{0.55em}
\setlist{itemsep=0.2em,topsep=0.3em,parsep=0.15em}
\titleformat{\section}{\Large\bfseries\sffamily\color{tinta}}{\thesection}{0.6em}{}[\vspace{-0.4em}{\color{tinta!35}\rule{\linewidth}{0.8pt}}]
\titleformat{\subsection}{\large\bfseries\sffamily\color{tinta}}{\thesubsection}{0.6em}{}
\titleformat{\subsubsection}{\normalsize\bfseries\sffamily\color{tinta2}}{\thesubsubsection}{0.6em}{}
\titlespacing*{\section}{0pt}{1.4em}{0.6em}
\titlespacing*{\subsection}{0pt}{1.0em}{0.3em}
\titlespacing*{\subsubsection}{0pt}{0.8em}{0.2em}

% ----------------------------------------------------------- macros
\newcommand{\s}[1]{\{#1\}}                         % conjunto por extensión
\newcommand{\vac}{\varnothing}
\newcommand{\cl}[1]{\overline{#1}}                 % clausura
\newcommand{\inte}[1]{\mathring{#1}}               % interior (letra sola)
\newcommand{\Int}{\operatorname{int}}
\newcommand{\Ext}{\operatorname{ext}}
\newcommand{\Fr}{\operatorname{Fr}}
\newcommand{\cotau}{\mathrm{co}\text{-}\tau}
\newcommand{\co}[1]{\mathrm{co}\text{-}#1}
\newcommand{\N}{\mathcal{N}}
\newcommand{\B}{\mathcal{B}}
\newcommand{\C}{\mathfrak{C}}
\newcommand{\Sb}{\mathcal{S}}
\newcommand{\Pt}{\mathcal{P}}
\newcommand{\R}{\mathbb{R}}
\newcommand{\Z}{\mathbb{Z}}
\newcommand{\Q}{\mathbb{Q}}
\newcommand{\sm}{\smallsetminus}
\newcommand{\homeo}{\cong}
\newcommand{\si}{\textcolor{verde}{\textbf{sí}}}
\newcommand{\no}{\textcolor{rojo}{\textbf{no}}}
\newcommand{\ok}{\textcolor{verde}{$\checkmark$}}
\newcommand{\mal}{\textcolor{rojo}{$\times$}}
\newcommand{\fuente}[1]{\textsuperscript{\textcolor{tinta2}{[#1]}}}

% \input expandible, para poder usarlo dentro de tablas antes de \bottomrule
\makeatletter\newcommand{\tinput}[1]{\@@input gen/#1.tex }\makeatother

% ------------------------------------------------------------ cajas
\tcbset{base/.style={enhanced,breakable,boxrule=0pt,arc=1.5pt,left=9pt,right=9pt,top=6pt,bottom=6pt,
  before skip=0.8em,after skip=0.8em,fonttitle=\bfseries\sffamily\small,coltitle=white,
  attach boxed title to top left={xshift=6pt,yshift=-2.5mm},boxed title style={arc=1.5pt,boxrule=0pt,top=1.5pt,bottom=1.5pt}}}
\newtcolorbox{definicion}[1][Definición]{base,colback=tinta!5,borderline west={2.5pt}{0pt}{tinta},
  colbacktitle=tinta,title={#1}}
\newtcolorbox{teorema}[1][Teorema]{base,colback=tinta2!6,borderline west={2.5pt}{0pt}{tinta2},
  colbacktitle=tinta2,title={#1}}
\newtcolorbox{ejemplo}[1][Ejemplo]{base,colback=verde!5,borderline west={2.5pt}{0pt}{verde},
  colbacktitle=verde,title={#1}}
\newtcolorbox{ojo}[1][Advertencia]{base,colback=ambar!7,borderline west={2.5pt}{0pt}{ambar},
  colbacktitle=ambar,title={#1}}
\newtcolorbox{metodo}[1][Método]{base,colback=papel,borderline west={2.5pt}{0pt}{black!55},
  colbacktitle=black!65,title={#1}}
\newtcolorbox{respuesta}{enhanced,boxrule=0.6pt,colframe=verde!70!black,colback=verde!4,arc=1.5pt,
  left=8pt,right=8pt,top=4pt,bottom=4pt,before skip=0.6em,after skip=0.8em}
\newcommand{\resp}[1]{\begin{respuesta}\textbf{\sffamily\color{verde!70!black}Respuesta.} #1\end{respuesta}}
\newtcolorbox{enunciado}{enhanced,breakable,boxrule=0pt,colback=black!4,arc=1.5pt,left=8pt,right=8pt,top=5pt,bottom=5pt,
  before skip=0.6em,after skip=0.6em,fontupper=\itshape}

% -------------------------------------------------- estilos de TikZ
\tikzset{
  hasse/.style={draw=tinta!45,line width=0.35pt},
  nodoh/.style={circle,draw=tinta,fill=white,inner sep=0.8pt,minimum size=5.6mm,font=\scriptsize},
  pto/.style={inner sep=2pt,font=\normalsize},
  flecha/.style={-{Stealth[length=2.2mm]},line width=0.55pt},
  ovalo/.style={draw=black!70,line width=0.5pt,ellipse,minimum width=2.1cm,minimum height=3.9cm},
}
% diagrama de flechas de una función entre conjuntos de 4 puntos:
% \diagrama{nombre}{d1,d2,d3,d4}{c1,c2,c3,c4}{i1/j1,i2/j2,...}
\newcommand{\diagrama}[4]{%
\begin{tikzpicture}[baseline=(current bounding box.center)]
  \node[ovalo] (L) at (0,0) {};
  \node[ovalo] (Rr) at (3.6,0) {};
  \node at (1.8,2.05) {$#1$};
  \foreach \t [count=\i] in {#2} {\node[pto] (l\i) at (0,1.5-\i*0.6) {$\t$};}
  \foreach \t [count=\i] in {#3} {\node[pto] (r\i) at (3.6,1.5-\i*0.6) {$\t$};}
  \foreach \a/\b in {#4} {\draw[flecha] (l\a) -- (r\b);}
\end{tikzpicture}}
